\subsection{The selected theta section in the incident frames}
\label{subsec:incident-sections}

The local jet estimate now requires identifying the theta section of the selected tuple in each incident frame.

Fix a word, one of its used twisting tuples and a selected coefficient
place \(w\) over \(p\). Use the determinant character, curve ideal,
twelfth root of the discriminant, translated theta vector and type
levels attached to this tuple. At consecutive ranks the parameters are
\[
 i=j-1,\quad g=t+1,\quad N=k+ig,\quad N'=N+g,
 \qquad A=ig/2,\quad A'=(i+1)g/2 .
\]
The ranks and \(g\) are even. Write \(\lambda\) for the embedding
of algebraic coefficients in \(\C_p\) extending the chosen place
on the coefficient field. All local absolute values below use this
embedding and \(v_p(p)=1\).

A label at level \(i\) consists of a component and a lattice \(L_v\)
at each split factor above \(p\).

\begin{proposition}[Companion lattices in the incident frames]\label{prop:incident-companion-lattices}
At two consecutive levels the label and its component frame give
self-dual rational companion lattices. Writing \(C,C'\) for the two companion lattices, their local spans and global self-duality are
\[
 C_v=F_vL_v,\qquad C'_v=F'_vL'_v,\qquad C^\vee=C,\qquad(C')^\vee=C'.
\]
The determinant ideles in the theta kernel are \(\det x,\det x'\).
\end{proposition}

\begin{proof}
To obtain its adelic point \(x\),
keep the component's away coordinates. At \(v\) use a matrix
carrying the reference lattice to \(L_v\) in the component's frame.
At the conjugate prime the unitary condition prescribes the inverse
adjoint matrix. Changing the chosen bases on the right changes the
point by the local level. Write 
\[
  x=\gamma y_b\kappa,
\]
 with \(y_b\)
the fixed representative in its level class, \(\gamma\) rational
unitary and \(\kappa\) in the compact level. At level zero use the
untwisted representative with the prescribed change at the twist
support. This change has trivial level character and is the identity
above \(p\). At positive levels the classes and representatives are
those of the actual step of the twisted chain.

\proofstep{Recover the local factors of the rational companion.}

Let \(M\) be the companion lattice of \(x\) in the rational orthonormal
space. If \(V\) is the rational orthonormal space and \(M_v\) are the translated local base lattices, the companion lattice is their rational intersection
\[
  M=\{u\in V:u_v\in M_v\text{ at every finite }v\}.
\]
 Localization recovers every
factor. Indeed, after multiplying by a common denominator these
are sublattices of one integral lattice with conditions modulo
finitely many relatively prime ideal powers. The Chinese remainder
theorem lifts each prescribed finite set of residues, to arbitrarily
high precision at a chosen prime.

For the integral lattice \(L\) after clearing denominators and coprime ideal powers \(\mathfrak a_v\), the residue map is
\[
  L/\bigl(\prod_v\mathfrak a_v\bigr)L\xrightarrow{\ \sim\ }\prod_v L/\mathfrak a_v L .
\]
 Dualizing locally commutes with
this construction. The standard companion is self-dual in the
chosen normalization, and the adelic action preserves the pairing.
Thus \(M\) is self-dual. At \(v\) its span is \(F_v L_v\), where
\(F_v\) takes the component frame into the orthonormal frame.
At the next level the same construction gives \(M'\) with span
\(F'_v L'_v\). The determinant ideles used for these labels are
\(\det x,\det x'\). This describes both the completed lattices and
the determinant entering the theta kernel. \label{cor:incident-sections-1}

\end{proof}

Write \(I,I'\) for their determinant ideals relative to the two
standard companions. The normalized theta vector on the pair is
the transported vacuum times
\[
 C_\chi=\chi(I)^{N'/2}\chi(I')^{-N/2}.
\]
\begin{proposition}[The normalized incident Taylor series]\label{prop:incident-taylor-series}
At a good inert coefficient place for the chosen tuple, the scalar \(S\) from \textup{\eqref{eq:tag-i1}} satisfies
\[
 v_p\!\left(p^{\sum_v\Delta_j(v)(l'_v-l_v)}S\right)=0.
\]
 In the
trace-dual pair frames the length-normalized series is \textup{\eqref{eq:tag-i4}},
with \(U\) integral and the same phase \(B_{d_v}\) in each entry
at \(v\). The factor \(U\) can depend on all places.
\end{proposition}

\begin{proof}
\proofstep{Normalize the finite sum of translated vacua.}

The finite reader \(\Xi\) weights its translates on the prescribed
quotient of the Hom lattice. Its coefficients are algebraic and
integral at \(w\). In fact the finite data are locally constant at
the stated level, and the denominators of the finite type projectors
are supported on the bad set for the tuple. Both primes at \(v\)
avoid this set. The translation quotient has level prime to \(p\).
The resulting kernel is a finite sum of entire translates with
algebraic Taylor coefficients in the chosen differential coordinates.

Write \(\mu\) for the stage character. Before the local length
normalization the scalar on this sum of normalized translates is
\[
 S=C_\chi^{-1}\frac{\mu(\det x')}{\mu(\det x)}
                  \chi(I)^{-A}\chi(I')^{A'} .
 \tag{I1}\label{eq:tag-i1}
\]
The determinant ideals of self-dual companions have norm one.
Their \(\chi\)-values have absolute value one at a good inert
prime. For example, a positive power of the ideal is principal,
and its prescribed character power is the reflex value of a
norm-one generator up to a root of unity. Complex conjugation
fixes the place on the quadratic reflex field and reverses this
value, so its valuation is zero.

\proofstep{Compute the valuation from the alternating cycle.}

For each \(v\), let \(s\) run around its cycle and write
\(\tau(v,s)\) for the corresponding infinite position. The
completion \(\rho_v\) of the host on the first factor, the
chosen inclusion \(\iota_v\) of its unramified coefficient field,
and the Witt Frobenius \(\sigma\) satisfy
\[
 \lambda(\tau(v,s)(z))=
 \iota_v\sigma^s\rho_v
 \begin{cases} z,&s\text{ even},\\ \bar z,&s\text{ odd}.
 \end{cases}
 \tag{I2}\label{eq:tag-i2}
\]
Indeed \(s=0\) fixes the prolongation and moving one place in
the real cycle composes with residue Frobenius. On the quadratic
CM factor this composition changes the structural embedding.
The CM-type lift therefore includes conjugation exactly at the
odd positions.

In particular a uniformizer idele \(\pi_v\) with
components \(p\) and \(p^{-1}\) at the conjugate places and one
elsewhere has \(v_p(\mu(\pi_v))=\Delta_j(v)\). The character's
infinity exponents have the stage weights after the common shift.
A constant shift disappears in the alternating sum on this even
cycle. One can compute this valuation on principal ideles after
raising to the exponent of the finite ray quotient, since compact
level factors have finite character. The local determinant order
of \(x\) is \(-l_v\) for log volume \(l_v\). Applying the formula
at every \(v\) gives
\[
    v_p(S)=-\sum_v\Delta_j(v)(l'_v-l_v).
 \tag{I3}\label{eq:tag-i3}
\]
Thus the factor \(p^{\sum_v\Delta_j(v)(l'_v-l_v)}\lambda(S)\)
in the local Taylor reading is a unit. The valuation statement
holds for lattices of either sign of length.

\proofstep{Pass from the orthogonal chart to the paired frames.}

Use one good model and one tangent unit \(c\) of the
formal isomorphism in Section~\ref{subsec:cm-models} for the
whole site. A basis of \(F_vL_v\) and its trace dual on the
conjugate summand identify the completed lattice with a
hyperbolic pair. Do this also for \(F'_vL'_v\) and take the
resulting Hom bases. In both the reference and these bases the
orthogonal model of the fourfold lattice has unimodular Gram
matrix. The transition matrix \(U\) between the two integral orthonormal bases satisfies
\[
  U^*U=I,\qquad U^{-1}=U^*\in M_n({\cal O}_w),
\]
where \(n\) is the size of the frame.

The change of frame \(U\) preserves the sum of the one-variable gauges modulo logarithms of integral principal units, by the
calculation of Section~\ref{subsec:cm-models}. The normalized
fourth tensor expansion of Section~\ref{subsec:theta-normalization}
applies to each translate, since the translation order is prime
to \(p\). Extracting the coefficient at four equal indices on
disjoint sets of variables proves integrality of the gauged
single section. Summing with the reader coefficients and using
\eqref{eq:tag-i3} preserves integrality. On each matrix entry the two
isotropic factors have zero self-pairings and identity cross
pairing. Substitution of the two trace-dual bases as in
Proposition~\ref{prop:curve-cycle-exponent} now gives the series
\[
    U\big((g_{d_v}(Z_{v,ab}))_{v,a,b}\big)
    \exp\left(-\sum_{v,a,b} B_{d_v}(Z_{v,ab})\right).
 \tag{I4}\label{eq:tag-i4}
\]
Here \(g_d\) is the vector of inverse logarithms on the cycle,
and \(U\) has integral coefficients in the completed algebraic
closure. All entries at \(v\) have the same \(B_{d_v}\).
The multiplier may involve variables from every place. The
series comparison does not require its coefficients to factor.

For a change of source and target bases \(C,C'\), the Hom
coordinate at the even positions transforms by
\(Z\mapsto C'ZC^{-1}\), at the odd positions by
\(Z\mapsto(C'^{-1})^t Z C^t\), with the corresponding
Frobenius applied to the matrices. Hence for an incident pair
one may take a basis adapted to the sublattice, with \(p\) in
the selected diagonal positions. This gives the scaling on the
two parities used in \eqref{eq:tag-j1}. \label{cor:incident-sections-2}

\end{proof}

\begin{proposition}[Equivariance and the local neighbor relation]\label{prop:incident-covariance}
The length-normalized incident series transforms under a rational
source move by \textup{\eqref{eq:tag-i5}} with its length correction and the
prescribed Hom substitution. Fix every coordinate away from \(v\). The character of the two length changes, with units \(a,a'\) from \textup{\eqref{eq:tag-i6}}, is
\[
 (n,n')\longmapsto a^n(a')^{n'}.
\]
 After removing
these units and the volume normalization, every rational indicator
relation of \(\{z:zL_v\subset L'_v\}\) holds coefficientwise for
the normalized series, with the remaining length factor of \textup{\eqref{eq:tag-p1}}.
\end{proposition}

\begin{proof}
\proofstep{Transform the source by a rational unitary map.}

Consider
an arithmetic move of the source by \(g_0\) in its rational
unitary group. Write \(x_1=g_0 x\kappa\) and \(e\) for \(g_0\)
conjugated into the orthonormal frame. Then \(M_1=eM\) and
\(F_{1,v}\gamma_v=\rho_v(e)F_v\), with \(\gamma_v\) the action
on the source coordinates. In the reader replace \(z\) by \(ze\).
The change of the finite character is
\(c_\Xi=\mu(\det\kappa)/\chi_{\rm lev}(\kappa)\).
Writing \(P=\prod_\tau \tau(\det g_0)\), the determinant ideal
rule gives \(\chi(I_1)=P^{-1}\chi(I)\).
Substitution in \eqref{eq:tag-i1} gives
\[
 \frac{S_1 c_\Xi}{S}
   =\frac{P^{A+N'/2}}{\mu(\det g_0)\chi_{\rm lev}(\kappa)} .
 \tag{I5}\label{eq:tag-i5}
\]
For the length-normalized series there is the additional factor
\(p^{-\sum_v\Delta_j(v)(l_{1,v}-l_v)}\).
This is exactly the inverse arithmetic character times the
inverse differential determinant, with its length correction.

To see the equality, the local differential determinant on
\(\gamma_v\) is the product of its embedded determinants raised
to \((-1)^s w_j(s)\) at positions \((v,s)\). Apply \eqref{eq:tag-i2} to
those factors, use the finite character on \(g_0 x\kappa\),
and use 
\[
  \operatorname{ord}_v\det\gamma_v=l_v-l_{1,v}
\]
in the length correction. The shifted constant exponent on
\(\det g_0\) gives \(A+N'/2\) in \eqref{eq:tag-i5}. With the Hom
substitution above this proves covariance on the whole orbit.
In particular two ways of carrying a representative to the
same lattice give equal kernel readings by applying covariance
to their quotient in the stabilizer. \label{cor:incident-sections-3}

\proofstep{Vary the two lattices at one split factor.}

Next fix every component and lattice away from \(v\), and vary
\(L_v,L'_v\) within this fiber. Set \(m_v=\lambda(\mu(\pi_v))\)
and let \(z_v\) be the \(\chi\)-value of its relative
determinant ideal. The units in \eqref{eq:tag-p1} are
\[
     a=p^{-\Delta_j(v)}m_v z_v^{A},
   \qquad a'=p^{\Delta_j(v)}m_v^{-1}z_v^{-A'} .
 \tag{I6}\label{eq:tag-i6}
\]
They belong to the embedded field of the readings. Indeed a
diagonal uniformizer in one split factor, completed by its dual,
realizes \(\pi_v\) as a determinant at either positive rank.
The step and twist characters on these adelic determinants take
their prescribed values in that field. The valuations follow from
\eqref{eq:tag-i2} and the norm-one ideal rule.

If the two lengths change by
\(n,n'\), the determinant ratio in the numerator of \eqref{eq:tag-i1},
including \(\chi(I)^{-A}\chi(I')^{A'}\), changes by
\(m_v^{n-n'}z_v^{An-A'n'}\). After multiplication by
\(p^{\Delta_j(v)(n'-n)}\) this is \(a^n(a')^{n'}\).

\proofstep{Transfer the indicator relation through the fixed away reader.}

Dividing the normalized theta vector by \(C_\chi\) restores the
transported vacuum. In the split model put \(H_v=\{z:zL_v\subset L'_v\}\). The two successive lattice changes give the normalized indicator
\[
  \operatorname{vol}(H_v)^{-1/2}\,\mathbf1_{H_v}(z).
\]
 This follows by composing the two changes of variables
in the Hom space, whose Jacobian gives the volume computed in
\eqref{eq:tag-p1}.

The trace different at these good places is a unit and
the reader away from \(v\) is fixed. Thus any rational linear
relation of indicators remains a relation of these normalized
series after removal of the two character units and the volume
factor. Indeed apply the identical away reader to the relation
before the finite sum of theta translates, then take each Taylor
coefficient. The length-normalized reading includes
\(p^{\Delta_j(v)(l'_v-l_v)}\). Removing it for the untwisted Hom
indicator calculation introduces precisely the last factor in
\eqref{eq:tag-p1}. This proves the required away-data statement for the full
neighbor sums. \label{cor:incident-sections-4}

\end{proof}

\begin{proposition}[The derivative and the neighbor operator use one scalar]\label{prop:incident-reading-scalar}
On the chosen stage and at the fixed coefficient place, the
class-representative derivative and the incident-frame derivative
are related by the scalar \(c_K\) of \textup{\eqref{eq:tag-i7}}. Writing \(v_w\) for the chosen valuation, the common scalar has the value
\[
 c_K=c^{-D}\lambda(k_T/(k_S c_{\rm op})),\qquad v_w(c_K)=0.
\]
It is independent of the lattice pair and works for every sharp polynomial of the given weight, also in the initial twisted copy. The separate length factor remains in the Taylor coefficient.
\end{proposition}

\begin{proof}
\proofstep{Pull back both class representatives.}

For decompositions
\(x=\gamma y_b\kappa_S\), \(x'=\gamma' y_{b'}\kappa_T\),
the pullback of the class kernel uses the rational unitary
maps on both factors. Its scalar on the theta function is
\[
     \frac{\mu(\det\kappa_S)\chi_{{\rm lev},T}(\kappa_T)}
          {\mu(\det\kappa_T)\chi_{{\rm lev},S}(\kappa_S)} .
\]
This follows by changing variables by the two rational maps in
the translation sum. Each finite weight transforms by its two
compact characters. Apply the chain rule to the frame minors
and contract their source indices against a polynomial of the
specified sharp weight. Each contracted index pairs a basis
matrix with its inverse. The residual determinant powers at
the source and target can be recorded as follows. For a rational
matrix \(M_0\) write \([M_0]_b=\prod_\tau\tau(\det M_0)^{b_\tau}\).

Let \(B_i\) be the chart basis in the rational orthonormal
coordinates and \(G_i\) the fixed change to orthonormal
coordinates for the identity element. Let \(q_i\) be the
level's determinant twist, and \(s_{\rm sh},s_{\rm div}\)
the source sharp and target divided shifts of the stage. Put
\[
 k_S=[B_i]_{-s_{\rm sh}}[G_i]_{N'/2+q_i},\qquad
 k_T=[B_{i+1}]_{-s_{\rm div}}[G_{i+1}]_{N/2+q_{i+1}} .
\]
Indeed the frame polynomial changes by the determinant to its
weight shift, and removing the determinant line of the theta
vector adds the constant half-rank to the twist exponent.
Let \(c_{\rm op}\) be the odd row reversal sign times the product
of the chosen derivation generators to minus the number of boxes
at their respective positions, and \(D\) the total number of boxes.
The differential comparison uses the single scalar
\[
                 c_K=c^{-D}\lambda(k_T/(k_S c_{\rm op})).
 \tag{I7}\label{eq:tag-i7}
\]
\proofstep{Identify the single differential scalar.}

For the cancellation, write \(G_i(\gamma)\) for the rational
change to orthonormal coordinates for \(\gamma\), so that
\(G_i(1)=G_i\). Put \(z_b=\chi_D(\det y_b)\) for the twist
character on the determinant class, and use the analogous
notation at the target. With \(P_i(\gamma)=[G_i(\gamma)]_1\)
the two residual factors are
\[
\begin{aligned}
 u_S&=\chi_{{\rm lev},S}(\kappa_S)^{-1}\mu(\det\gamma)^{-1}
            P_i(\gamma)^{N'/2+q_i} z_b^{-q_i},\\
 u_T&=\chi_{{\rm lev},T}(\kappa_T)\mu(\det\gamma')
            P_{i+1}(\gamma')^{-(N/2+q_{i+1})} z_{b'}^{q_{i+1}} .
\end{aligned}
\]

\proofstep{Cancel the level characters and rational determinant factors.}

To obtain them, insert both double-coset decompositions in
\eqref{eq:tag-i1} times the compact character ratio. The determinant
characters of the representatives give the two class characters
already in the stage differential. Moving the rational factors
out of the determinant ideals by \eqref{eq:tag-i5} and its target version
leaves \(u_Su_T\). The separate power of \(p\) given by the
lengths stays in the local Taylor reading. Include the source
coordinate determinant \(c_A\) with the signed weights of \eqref{eq:tag-i5},
and write \(s_O\) for the source reading with its level character,
twist, inverse basis determinant to the twist and operator
normalization. Multiplication by \(u_S\) cancels the source
level character and \(z_b^{q_i}\). The rational and frame
determinants combine into the two factors of \(k_S\). Thus

\[
  c_A s_O\lambda(u_S)=\lambda(k_S c_{\rm op}).
\]
 The target
reading equals \(\lambda(u_T k_T)c^{-D}\), since each box
differentiates the argument scaled by \(c\). Consequently
\[
  c_K(c_A s_O)\lambda(u_S)\lambda(u_T)=\lambda(u_T k_T)c^{-D}.
\]
This proves the scalar identity with \eqref{eq:tag-i7}.
This computation applies coefficientwise to all the sharp
polynomials, including those in the initial twisted copy of
level zero, whose frames are identified by its level isomorphism.

\proofstep{Check the valuation and independence from the pair.}

The bases \(B_i\) are shared across components of a stage in the
chosen chart construction. Their embedded determinants have norm
one. Locally they compare integral bases on the given split
lattice and its dual with the integral bases on the standard
companion, so the local change is invertible over the valuation
ring in the relevant embedding. For the constant powers of \(G_i\)
one needs only the product over the CM type. Its determinant is
norm one and the reflex product is in the quadratic CM field,
where the inert valuation is invariant under conjugation.

For the norm-one reflex product \(u\) at this inert place the calculation is
\[
  v_p(u)=v_p(\bar u),\qquad 2v_p(u)=v_p(u\bar u)=0 .
\]
 The
product therefore has valuation zero. The derivation generators
are units at this good place, as is \(c\). These facts give
\(|c_K|_p=1\) and show that \eqref{eq:tag-i7} is independent of the lattice
pair. Consequently the calculation of the neighbor operator
before differentiation and the application of \eqref{eq:tag-j1} after
differentiation give bounds on the same stage of the selected
chain. This common scalar allows the linear combinations in the
Euler inversion of \eqref{eq:tag-p7}. \label{cor:incident-sections-5}
\end{proof}
